\documentclass[11pt]{article}
\usepackage[letterpaper,margin=1in]{geometry}
\usepackage[T1]{fontenc}
\usepackage{lmodern}
\usepackage{microtype}
\usepackage{amsmath,amssymb,amsthm,mathtools}
\usepackage{xcolor}
\usepackage{hyperref}
\hypersetup{colorlinks=true,linkcolor=blue!45!black,
 citecolor=blue!45!black,urlcolor=blue!45!black,
 pdftitle={Optimal Berry--Esseen bounds for bounded-batch discovery},
 pdfauthor={Christopher D. Long}}
\usepackage{fancyhdr}
\pagestyle{fancy}
\fancyhf{}
\fancyhead[L]{\small Optimal bounds for bounded-batch discovery}
\fancyhead[R]{\small Research draft}
\fancyfoot[C]{\thepage}
\setlength{\headheight}{14pt}
\setlength{\emergencystretch}{2em}
\numberwithin{equation}{section}
\newtheorem{theorem}{Theorem}[section]
\newtheorem{lemma}[theorem]{Lemma}
\newtheorem{proposition}[theorem]{Proposition}
\newtheorem{corollary}[theorem]{Corollary}
\theoremstyle{remark}
\newtheorem{remark}[theorem]{Remark}
\newcommand{\E}{\mathbb E}
\newcommand{\Pp}{\mathbb P}
\newcommand{\N}{\mathbb N}
\newcommand{\R}{\mathbb R}
\newcommand{\C}{\mathbb C}
\newcommand{\ind}{\mathbf 1}
\newcommand{\Var}{\operatorname{Var}}
\newcommand{\Poi}{\operatorname{Poi}}
\newcommand{\Bin}{\operatorname{Bin}}
\newcommand{\dK}{d_{\mathrm K}}
\newcommand{\Normal}{\mathcal N(0,1)}
\newcommand{\calZ}{\mathcal Z}
\title{Optimal Berry--Esseen bounds for bounded-batch discovery}
\author{Christopher D. Long}
\date{November 26, 2026\\[3pt]\small Research draft}

\begin{document}
\maketitle

\begin{abstract}
Independent random finite batches arrive according to a Poisson process, and
$X$ counts the species discovered by a fixed time. If every batch has at most
$B$ species and $V=\Var(X)>0$, we prove
$\dK((X-\E X)/\sqrt V,\Normal)\le\min\{1,150B/\sqrt V\}$.
The constant is uniform in the batch law, time, and countable species space;
no bound on a species dependence degree is imposed. A bipartite
independence-polynomial representation gives a zero-free disk of radius
proportional to $B^{-1}$ for the moment generating function. Combining this
with the actual-variance moment inequality proved in the companion paper
gives variance-linear cumulant bounds and a Fourier smoothing proof of the
Berry--Esseen estimate. The same argument applies to nonnegative weighted
discovery with total weight at most $W$ in each batch, yielding
$150W/\sqrt{V_w}$ even without a cardinality bound. Disjoint batches show that
the order $B/\sqrt V$ is optimal. The statements concern scalar Poissonized
counts; quantitative fixed-sample transfer and higher-coverage counts require
additional arguments.
\end{abstract}

\section{Model and main results}\label{sec:main}

Let $\mathcal S$ be a countable set. Let $A_1,A_2,\ldots$ be independent,
identically distributed random finite subsets of $\mathcal S$, and let
$(N(t))_{t\ge0}$ be an independent rate-one Poisson process. Empty batches are
allowed. Define the discovery set and its cardinality by
\[
 \mathcal D(t)=\bigcup_{k=1}^{N(t)}A_k,
 \qquad X(t)=|\mathcal D(t)|.
\]
For real random variables $U$ and $Z$, write
\[
 \dK(U,Z)=\sup_{x\in\R}|\Pp(U\le x)-\Pp(Z\le x)|.
\]
We use the same notation with a probability distribution in the second
argument. Throughout, $\Normal$ denotes the standard normal law.

\begin{theorem}[Bounded batch cardinality]\label{thm:unweighted}
Suppose $|A_1|\le B$ almost surely, where $B\in\N$. For every $t\ge0$ such
that $V(t)=\Var(X(t))>0$,
\begin{equation}\label{eq:mainBE}
 \dK\left(\frac{X(t)-\E X(t)}{\sqrt{V(t)}},\Normal\right)
 \le\min\left\{1,\frac{150B}{\sqrt{V(t)}}\right\}.
\end{equation}
The constant is independent of the batch law, time, number of species, and
overlap pattern.
\end{theorem}

The theorem includes sequences in which all these parameters vary. In
particular, $B^2/V\to0$ implies asymptotic normality. It also covers arbitrarily
dense overlap: a single species can share positive-probability batches with
arbitrarily many other species.

Here is the more general statement proved below. Fix weights $h_i\ge0$ for
$i\in\mathcal S$ and set
\[
 X_h(t)=\sum_{i\in\mathcal D(t)}h_i,
 \qquad V_h(t)=\Var(X_h(t)).
\]

\begin{theorem}[Bounded total weight]\label{thm:weighted}
Suppose that, for some $W>0$,
\begin{equation}\label{eq:weightbound}
 \sum_{i\in A_1}h_i\le W\qquad\text{almost surely}.
\end{equation}
Batch cardinalities need not be uniformly bounded. For every $t\ge0$ with
$V_h(t)>0$,
\begin{equation}\label{eq:weightedBE}
 \dK\left(\frac{X_h(t)-\E X_h(t)}{\sqrt{V_h(t)}},\Normal\right)
 \le\min\left\{1,\frac{150W}{\sqrt{V_h(t)}}\right\}.
\end{equation}
Moreover, for every integer $k\ge3$, the cumulants satisfy
\begin{equation}\label{eq:maincumulants}
 |\kappa_k(X_h(t))|\le2k!V_h(t)(4W)^{k-2}.
\end{equation}
\end{theorem}

Theorem~\ref{thm:unweighted} follows by taking $h_i=1$ and $W=B$.
The bound $X_h(t)\le WN(t)$ ensures all moments and exponential moments needed
below. Species of weight zero, and species that are never sampled, may be
discarded without changing $X_h$.

\subsection*{The companion-paper input and the new step}

The input from the companion paper is an actual-variance estimate on the real
axis. In the notation above, with $M=X_h(t)-\E X_h(t)$ and $V=V_h(t)$, it is
\begin{equation}\label{eq:realmgf}
 \log\E e^{\theta M}
 \le\frac{V\theta^2}{2(1-2W|\theta|)},
 \qquad \theta\in\R,\quad |\theta|<\frac1{2W}.
\end{equation}
For unweighted discovery this is
\cite[Proposition~9.1]{LongCompanion}; its weighted extension, including
unbounded batch cardinalities under \eqref{eq:weightbound}, is
\cite[Theorem~11.1]{LongCompanion}. These are explicit imported results,
with their full proofs in the companion manuscript. Section~\ref{sec:input}
recalls the probabilistic ingredients to identify precisely what is used.

We prove that the complex moment generating function is zero-free for
$|z|<1/(2W)$. The real-axis bound \eqref{eq:realmgf} then bounds the real part
of its analytic logarithm by a constant times $V/W^2$ on a smaller disk.
A harmonic coefficient estimate gives \eqref{eq:maincumulants}, and Esseen
smoothing gives \eqref{eq:weightedBE}. A zero-free disk alone does not supply
the variance-linear coefficient estimate used in this argument. The
actual-variance input is what makes the final error depend on $W/\sqrt V$
instead of the number of species. The general connection between zeros and
normal approximation is developed, for example, in \cite{LPRS}; the proof
here gives the required coefficients directly.

\section{A bipartite zero-free lemma}\label{sec:bipartite}

For a finite graph $G$ with vertex activities $(\lambda_v)$, let
\[
 \calZ_G(\boldsymbol\lambda)
 =\sum_{I\subseteq V(G):\ I\text{ independent}}\prod_{v\in I}\lambda_v
\]
be its multivariate independence polynomial, with
$\calZ_\varnothing=1$. All induced subgraphs retain their original vertex
activities.

\begin{lemma}[Complex activities on one side]\label{lem:bipartite}
Let $G=(L\cup R,E)$ be a finite bipartite graph. Assign activities
$\lambda_i=z_i\in\C$ for $i\in L$ and $\lambda_j=w_j\ge0$ for $j\in R$.
Suppose that numbers $0\le r_i<1$, $i\in L$, satisfy
\begin{equation}\label{eq:arcsincondition}
 |z_i|\le r_i,
 \qquad \sum_{i\in N_G(j)}\arcsin r_i<\frac\pi2
 \quad(j\in R).
\end{equation}
Then $\calZ_G\ne0$. More precisely, for every induced subgraph $H$ and every
$v\in V(H)$, the ratio
\begin{equation}\label{eq:ratio}
 R_v(H)=\lambda_v
   \frac{\calZ_{H-v-N_H(v)}}{\calZ_{H-v}}
\end{equation}
is defined and satisfies
\begin{equation}\label{eq:ratiobounds}
 |R_i(H)|\le r_i\quad(i\in L\cap V(H)),
 \qquad \operatorname{Re}R_j(H)\ge0\quad(j\in R\cap V(H)).
\end{equation}
\end{lemma}

\begin{proof}
We prove nonvanishing and \eqref{eq:ratiobounds} simultaneously by induction
on $|V(H)|$. The empty partition function is one. Assume all assertions hold
for proper induced subgraphs of a nonempty $H$. Every denominator in
\eqref{eq:ratio} is therefore nonzero.

Fix $v$ and enumerate its neighbors in $H$ as $v_1,\ldots,v_d$. Define
\[
 H_\ell=H-v-\{v_1,\ldots,v_{\ell-1}\},\qquad 1\le\ell\le d.
\]
Partitioning independent sets according to whether they contain a specified
vertex gives
\begin{equation}\label{eq:deleteidentity}
 \calZ_H=\calZ_{H-v}(1+R_v(H)).
\end{equation}
Apply this identity successively in the smaller graphs $H_\ell$ to obtain
the exact recursion
\begin{equation}\label{eq:recursion}
 R_v(H)=\lambda_v\prod_{\ell=1}^d
                  (1+R_{v_\ell}(H_\ell))^{-1}.
\end{equation}
Every ratio on the right belongs to a proper induced subgraph, and all its
denominators are nonzero by induction.

If $v=i\in L$, its neighbors are in $R$. The induction hypothesis gives
$\operatorname{Re}R_{v_\ell}(H_\ell)\ge0$, hence
$|1+R_{v_\ell}(H_\ell)|\ge1$. Equation~\eqref{eq:recursion} yields
$|R_i(H)|\le|z_i|\le r_i$.

If $v=j\in R$, its neighbors are in $L$. For $|\zeta|\le r<1$, the disk
centered at one of radius $r$ lies in the right half-plane, and its two
tangent rays from the origin give
\begin{equation}\label{eq:argdisk}
 |\arg(1+\zeta)|\le\arcsin r.
\end{equation}
If $w_j=0$, then $R_j(H)=0$. Otherwise, use the principal arguments of the
factors in \eqref{eq:recursion}; their sum has absolute value at most
\[
 \sum_{i\in N_H(j)}\arcsin r_i
 \le\sum_{i\in N_G(j)}\arcsin r_i<\frac\pi2.
\]
Since $w_j$ is positive real, \eqref{eq:recursion} gives
$\operatorname{Re}R_j(H)>0$. This proves \eqref{eq:ratiobounds}.

For a left vertex, $|R_v(H)|<1$; for a right vertex,
$\operatorname{Re}(1+R_v(H))\ge1$. Thus $1+R_v(H)\ne0$ in either case.
Equation~\eqref{eq:deleteidentity} and the induction hypothesis now give
$\calZ_H\ne0$, completing the induction.
\end{proof}

\begin{corollary}[One-sided degree bound]\label{cor:bipartite}
If every vertex in $R$ has degree at most $B\ge1$, then $\calZ_G\ne0$
whenever $w_j\ge0$ for $j\in R$ and
\begin{equation}\label{eq:uniformradius}
 |z_i|<\sin\frac\pi{2B}\qquad(i\in L).
\end{equation}
No restriction is imposed on degrees in $L$.
\end{corollary}
\begin{proof}
Choose $r$ strictly between $\max_{i\in L}|z_i|$ and
$\sin(\pi/(2B))$, and apply Lemma~\ref{lem:bipartite} with $r_i=r$.
If $L$ is empty, the polynomial is a product of positive numbers $1+w_j$.
\end{proof}

\section{The absence polynomial and the moment generating function}
\label{sec:absence}

Fix $t\ge0$. First restrict discovery to a finite species set $J$.
Intersect each batch with $J$, discard empty intersections, and coalesce
identical remaining types. Independent thinning of a marked Poisson process
gives independent Poisson counts for these types. Retain only types of
positive sampling probability. For a nonempty type
$C\subseteq J$, write its mean as $a_C\ge0$ and set
\[
 q_C=e^{-a_C}>0,\qquad w_C=e^{a_C}-1\ge0.
\]
There are finitely many such types. Let
$Y_i=\ind_{\{i\notin\mathcal D(t)\}}$ for $i\in J$. The multivariate
absence polynomial is
\begin{align}
 Z_J(\boldsymbol u)
 &=\E\prod_{i\in J}(1+u_iY_i)\notag\\
 &=\sum_{S\subseteq J}\left(\prod_{i\in S}u_i\right)
                  \prod_{C:C\cap S\ne\varnothing}q_C.
 \label{eq:absencepoly}
\end{align}
The coefficient formula holds because all species in $S$ are absent exactly
when every batch type meeting $S$ has Poisson count zero.

Construct the bipartite incidence graph with left vertices $i\in J$, right
vertices the nonempty batch types $C$, and edges $i\sim C$ if $i\in C$.
Assign activity $u_i$ to a species vertex and activity $w_C$ to a batch
vertex. Conditional on the set $S$ of left vertices in an independent set,
the permitted right vertices are precisely those disjoint from $S$.
Consequently
\[
 \calZ_G
 =\sum_{S\subseteq J}\left(\prod_{i\in S}u_i\right)
                  \prod_{C:C\cap S=\varnothing}(1+w_C).
\]
Since $1+w_C=q_C^{-1}$, we have the exact identity
\begin{equation}\label{eq:gasidentity}
 Z_J(\boldsymbol u)=\left(\prod_Cq_C\right)\calZ_G.
\end{equation}
The prefactor is strictly positive. Isolated species vertices are harmless:
they contribute the factor $1+u_i$.

\begin{proposition}[Finite-species zero-free disks]\label{prop:finitezero}
For unweighted discovery with batches of size at most $B$,
\begin{equation}\label{eq:absencezerofree}
 Z_J(\boldsymbol u)\ne0
 \quad\text{whenever}\quad |u_i|<\sin\frac\pi{2B}\quad(i\in J).
\end{equation}
In particular, for $X_J=|\mathcal D(t)\cap J|$,
\begin{equation}\label{eq:unweightedzerofree}
 \E e^{zX_J}\ne0
 \qquad\text{if}\qquad
 |z|<\log\left(1+\sin\frac\pi{2B}\right).
\end{equation}
Under the weighted assumption \eqref{eq:weightbound}, let
$X_{h,J}=\sum_{i\in\mathcal D(t)\cap J}h_i$. Then
\begin{equation}\label{eq:weightedzerofree}
 \E e^{zX_{h,J}}\ne0\qquad\text{if}\qquad |z|<\frac1{2W}.
\end{equation}
\end{proposition}

\begin{proof}
Each batch vertex in the unweighted incidence graph has degree at most $B$.
Corollary~\ref{cor:bipartite} and \eqref{eq:gasidentity} give
\eqref{eq:absencezerofree}. Since $X_J=|J|-\sum_{i\in J}Y_i$,
\begin{equation}\label{eq:mgfabsence}
 \E e^{zX_J}=e^{z|J|}Z_J(e^{-z}-1,\ldots,e^{-z}-1).
\end{equation}
The estimate $|e^{-z}-1|\le e^{|z|}-1$ proves
\eqref{eq:unweightedzerofree}.

For the weighted assertion, discard zero-weight species and species that
never appear. Then $0<h_i\le W$ for each remaining species. The weighted
version of \eqref{eq:mgfabsence} is
\begin{equation}\label{eq:weightedmgfabsence}
 \E e^{zX_{h,J}}
 =e^{z\sum_{i\in J}h_i}
       Z_J\bigl((e^{-zh_i}-1)_{i\in J}\bigr).
\end{equation}
Fix $|z|<1/(2W)$ and put $r_i=e^{|z|h_i}-1$. For $0\le x\le1/2$,
$e^x-1\le2x$, so $0\le r_i<1$ and, for every restricted batch type $C$,
\[
 \sum_{i\in C}r_i\le2|z|\sum_{i\in C}h_i\le2|z|W<1.
\]
Convexity of $\arcsin$ on $[0,1]$ gives
$\arcsin r\le(\pi/2)r$. Hence
\[
 \sum_{i\in C}\arcsin r_i<\frac\pi2.
\]
Also $|e^{-zh_i}-1|\le r_i$. Apply Lemma~\ref{lem:bipartite} to
\eqref{eq:gasidentity}, and then \eqref{eq:weightedmgfabsence}, to obtain
\eqref{eq:weightedzerofree}. This part of the argument uses total batch
weight, with no bound on batch cardinality.
\end{proof}

\begin{proposition}[Countable species]\label{prop:countable}
Under \eqref{eq:weightbound}, the moment generating function of $X_h(t)$ is
entire and is zero-free for $|z|<1/(2W)$. The same is true for its centered
moment generating function.
\end{proposition}
\begin{proof}
Take an increasing finite exhaustion $J_n$ of the remaining species. Put
$F_n=X_{h,J_n}$ and $F=X_h(t)$. Then $0\le F_n\uparrow F\le WN(t)$ almost
surely. For every finite $A>0$,
\[
 \sup_{|z|\le A}|e^{zF_n}-e^{zF}|
 \le A(F-F_n)e^{AF}\le AF e^{AF}.
\]
The last bound is integrable, and the middle term tends to zero. Thus
$\E e^{zF_n}\to\E e^{zF}$ locally uniformly on $\C$, and the limiting
function is entire. By Proposition~\ref{prop:finitezero}, every function in
the sequence is zero-free on the disk $|z|<1/(2W)$. Hurwitz's theorem gives
that their limit is either identically zero or zero-free on that connected
disk. Its value at zero is one, excluding the first possibility. Centering
multiplies the function by the nonvanishing factor $e^{-z\E F}$.
\end{proof}

\section{The actual-variance input}\label{sec:input}

For clarity, we recall how the companion results cited in
\eqref{eq:realmgf} apply to this model. For fixed terminal time $t$, let
$M_s=\E[X_h(t)\mid\mathcal F_s]-\E X_h(t)$ be the discovery martingale.
If species $i$ has marginal arrival rate $\lambda_i$, its jump when a batch
$C$ arrives at time $s$ contributes
\[
 H_C(s)=\sum_{i\in C}h_i e^{-\lambda_i(t-s)}
                    \ind_{\{i\text{ absent just before }s\}},
 \qquad 0\le H_C(s)\le W.
\]
The companion paper gives a positive weighted-overlap decomposition of its
predictable quadratic variation $Q_t$ with coefficient load at most $2W^2$.
In particular, writing $V=\Var(X_h(t))$ and $\psi(u)=e^u-1-u$,
\begin{equation}\label{eq:bracketinput}
 \E Q_t=V,\qquad
 \log\E e^{a(Q_t-V)}\le\frac{V}{2W^2}\psi(2W^2a),
 \quad a\in\R.
\end{equation}
These are part of \cite[Theorem~11.1]{LongCompanion}. The weighted
incidence bound used there is
\[
 \sum_jh_j\lambda_{ij}
 =\E\left[\ind_{\{i\in A_1\}}\sum_{j\in A_1}h_j\right]
 \le W\lambda_i,
\]
so no bounded-cardinality hypothesis is needed for that result.

For reference, the short moment argument from
\cite[Proposition~9.1]{LongCompanion} uses the compensated-Poisson
exponential martingale and H\"older's inequality. With
$\psi_W(u)=(e^{Wu}-1-Wu)/W^2$ for $u\ge0$, it gives, for $q>1$,
\begin{equation}\label{eq:holderinput}
 \log\E e^{\theta M_t}
 \le\frac{q-1}{q}\frac{V}{2W^2}
 \left[\exp\left\{\frac{2W^2\psi_W(q|\theta|)}{q-1}\right\}-1\right].
\end{equation}
To check the stated denominator, put $b=W|\theta|\in(0,1/2)$ and $q=1+b$.
The series estimates
\[
 e^u-1-u\le\frac{u^2}{2(1-u/3)}\quad(0\le u<3),
 \qquad e^v-1\le\frac{v}{1-v/2}\quad(0\le v<2)
\]
bound \eqref{eq:holderinput} by
\[
 \frac{V\theta^2}{2}\,
 \frac{q}{1-qb/3-q^2b/2}.
\]
All denominators are positive, since
\[
 1-\frac{qb}{3}-\frac{q^2b}{2}-q(1-2b)
 =\frac b6+\frac{2b^2}{3}-\frac{b^3}{2}\ge0.
\]
This proves \eqref{eq:realmgf} from the imported bracket estimate and its
standard exponential-martingale consequence. At $\theta=0$ the estimate is
immediate. We will use only \eqref{eq:realmgf} in the remaining analysis.

\section{Variance-linear cumulants and Fourier smoothing}\label{sec:analytic}

The analytic step is useful independently of the discovery model.

\begin{lemma}[From a zero-free disk and a real moment bound]
\label{lem:analytic}
Let $M$ be a centered real random variable of variance $V>0$. Suppose its
moment generating function is analytic and zero-free for $|z|<1/(2L)$,
where $L>0$, and
\begin{equation}\label{eq:generalmgf}
 \log\E e^{\theta M}
 \le\frac{V\theta^2}{2(1-2L|\theta|)}
 \qquad(\theta\in\R,\ |\theta|<1/(2L)).
\end{equation}
Then, for every $k\ge3$,
\begin{equation}\label{eq:cumulantbound}
 |\kappa_k(M)|\le2k!V(4L)^{k-2},
\end{equation}
and
\begin{equation}\label{eq:abstractBE}
 \dK(M/\sqrt V,\Normal)
 \le\min\left\{1,C_*\frac L{\sqrt V}\right\},
 \qquad
 C_*:=\frac{36}{\sqrt\pi}+\frac{864}{\pi\sqrt{2\pi}}<131.
\end{equation}
\end{lemma}

\begin{proof}
The zero-free disk is simply connected, so there is an analytic logarithm
\[
 K(z)=\log\E e^{zM},\qquad K(0)=0.
\]
It satisfies $K'(0)=0$ and $K''(0)=V$. Set $R=1/(4L)$. For $|z|\le R$,
the modulus inequality and \eqref{eq:generalmgf} give
\begin{align}
 \operatorname{Re}K(z)
 &=\log|\E e^{zM}|\notag\\
 &\le\log\E e^{(\operatorname{Re}z)M}
 \le\frac{V(\operatorname{Re}z)^2}{2(1-2L|\operatorname{Re}z|)}
 \le VR^2.
 \label{eq:realpartbound}
\end{align}

Write $K(z)=\sum_{k\ge1}c_kz^k$. The function
$H(\vartheta)=VR^2-\operatorname{Re}K(Re^{i\vartheta})$ is nonnegative
and has mean $VR^2$, by the mean-value property and $K(0)=0$. For $k\ge1$,
Fourier coefficient extraction gives
\[
 c_kR^k=-\frac1\pi\int_0^{2\pi}
                      H(\vartheta)e^{-ik\vartheta}\,d\vartheta,
 \qquad |c_k|R^k\le2VR^2.
\]
This is the elementary harmonic proof of the Carath\'eodory coefficient
estimate. Since $c_k=\kappa_k(M)/k!$, it proves
\eqref{eq:cumulantbound}. Summing the bounds for $k\ge3$ yields, for
$|z|<R$,
\begin{align}
 \left|K(z)-\frac{Vz^2}{2}\right|
 &\le\sum_{k\ge3}2VR^2\left(\frac{|z|}{R}\right)^k\notag\\
 &=\frac{2V|z|^3}{R(1-|z|/R)}
 =\frac{8LV|z|^3}{1-4L|z|}.
 \label{eq:remainder}
\end{align}

Put $\sigma=\sqrt V$, $\delta=L/\sigma$, and $U=M/\sigma$. For
$|s|<1/(4\delta)$, equation~\eqref{eq:remainder} gives
\[
 K(is/\sigma)=-s^2/2+\rho(s),
 \qquad |\rho(s)|\le\frac{8\delta|s|^3}{1-4\delta|s|}.
\]
Take $T=1/(36\delta)$. For $|s|\le T$, the denominator in the last
display is at least $8/9$, so $|\rho(s)|\le9\delta|s|^3\le s^2/4$.
Using $|e^v-1|\le|v|e^{|v|}$, the characteristic function $\phi_U$ satisfies
\begin{equation}\label{eq:cfbound}
 |\phi_U(s)-e^{-s^2/2}|
 \le9\delta|s|^3e^{-s^2/4},\qquad |s|\le T.
\end{equation}

The normal-reference form of Esseen's smoothing inequality is
\begin{equation}\label{eq:Esseen}
 \dK(U,\Normal)
 \le\frac1\pi\int_{-T}^{T}
       \frac{|\phi_U(s)-e^{-s^2/2}|}{|s|}\,ds
       +\frac{24}{\pi\sqrt{2\pi}\,T}.
\end{equation}
This precise form appears in \cite[equation~(2.16)]{LPRS}, which cites
\cite[Section~XVI.4, Lemma~2]{Feller}. Since
\[
 \int_{-\infty}^{\infty}s^2e^{-s^2/4}\,ds=4\sqrt\pi,
\]
substitution of \eqref{eq:cfbound} into \eqref{eq:Esseen} gives
\[
 \dK(U,\Normal)
 \le\left(\frac{36}{\sqrt\pi}
        +\frac{864}{\pi\sqrt{2\pi}}\right)\delta.
\]
The constant in parentheses is $130.027828\ldots<131$. The Kolmogorov
distance is also at most one, proving \eqref{eq:abstractBE}.
\end{proof}

\begin{proof}[Proof of Theorem~\ref{thm:weighted}]
Proposition~\ref{prop:countable} gives a zero-free disk $|z|<1/(2W)$ for
the centered moment generating function. The real bound
\eqref{eq:realmgf} is supplied by the companion results as explained in
Section~\ref{sec:input}. Apply Lemma~\ref{lem:analytic} with $L=W$ and
$V=V_h(t)$. Cumulants of order at least two are unchanged by centering,
which gives \eqref{eq:maincumulants}; since $C_*<150$, it also gives
\eqref{eq:weightedBE}.
\end{proof}

\section{Optimal order and scope}\label{sec:scope}

\begin{proposition}[A matching order obstruction]\label{prop:sharpness}
The factor $B/\sqrt V$ in Theorem~\ref{thm:unweighted} cannot be replaced
uniformly by a quantity of smaller order as $V/B^2\to\infty$.
\end{proposition}
\begin{proof}
Fix an integer $B\ge1$. For each integer $m$, take $m$ disjoint batches of cardinality
$B$, assigning probability $1/m$ to each. At time $t_m=m\log2$, their
independent Poisson counts have mean $\log2$, so the batch-presence
indicators are independent Bernoulli variables of parameter $1/2$. Thus
\[
 X(t_m)\ \stackrel{d}{=}\ B\,\Bin(m,1/2),
 \qquad V(t_m)=\frac{B^2m}{4}.
\]
If a distribution has an atom of mass $p$ and the comparison distribution
is continuous, their Kolmogorov distance is at least $p/2$: the two limits
of the distribution-function difference on either side of the jump differ
by $p$. The largest binomial atom is
\[
 p_m=2^{-m}\binom{m}{\lfloor m/2\rfloor}
 \sim\sqrt{\frac2{\pi m}}.
\]
Consequently
\[
 \liminf_{m\to\infty}\frac{\sqrt{V(t_m)}}{B}
 \dK\left(\frac{X(t_m)-\E X(t_m)}{\sqrt{V(t_m)}},\Normal\right)
 \ge\frac1{2\sqrt{2\pi}}>0.
\]
The same obstruction applies to the weighted theorem by using disjoint
singletons each of weight $W$.
\end{proof}

The word ``optimal'' in the title concerns this order, not the numerical
constant 150. The proof uses substantial slack in its Taylor cutoff and
smoothing estimates.

The statements are scalar and Poissonized. For a deterministic number $n$
of batches, the independent Poisson counts used in
\eqref{eq:absencepoly} are unavailable. The companion paper proves a
fixed-law central limit theorem by de-Poissonization, but its asymptotic
transfer does not by itself give a uniform fixed-sample Berry--Esseen bound
of the form \eqref{eq:mainBE}; quantitative transfer errors and the
fixed-sample variance must also be controlled.

Likewise, counts of species observed at least $r\ge2$ times do not have the
absence-polynomial representation \eqref{eq:absencepoly}. The new
zero-free argument therefore does not automatically extend the optimal rate
to those counts. The argument uses nonnegative weights throughout; no
signed-weight version is asserted.

The result is an approximation of distribution functions. It does not
establish a span-one local limit theorem, and the disjoint-batch example
exhibits the relevant lattice obstruction. Relative Gaussian-tail
approximations require further analysis of the cumulants and of the
applicable range. We make no exhaustive priority claim for the elementary
bipartite lemma or the analytic tools used here.

\begin{thebibliography}{9}

\bibitem{Feller}
W. Feller.
\emph{An Introduction to Probability Theory and Its Applications},
Vol.~II, 2nd ed. John Wiley \& Sons, New York, 1971.

\bibitem{LPRS}
J. L. Lebowitz, B. Pittel, D. Ruelle, and E. R. Speer.
Central limit theorems, Lee--Yang zeros, and graph-counting polynomials.
\emph{Journal of Combinatorial Theory, Series A} \textbf{141} (2016),
147--183.
\href{https://doi.org/10.1016/j.jcta.2016.02.009}{doi:10.1016/j.jcta.2016.02.009};
\href{https://arxiv.org/abs/1408.4153}{arXiv:1408.4153}.
The smoothing inequality is equation~(2.16) in both the published paper and
the second arXiv version.

\bibitem{LongCompanion}
C. D. Long.
\emph{A fluctuation theory for bounded-batch species discovery}.
Research draft, November 26, 2026.
Proposition~9.1 and Theorem~11.1.
\href{https://github.com/long-mathematics/bounded-batch-discovery}{Repository},
\path{papers/bounded_batch_discovery.tex}.

\end{thebibliography}
\end{document}
